\documentclass[12pt,a4paper]{article}
        \makeatletter
        \def\input@path{{../}}
        \makeatother
        \usepackage{almeria-fichas}

        \FichaMeta{Sesion 09}{Perimetro y area: no son lo mismo}{Bloque 2 · Geometria urbana}{Comparar perimetro y area para no confundir ambas medidas al analizar zonas urbanas.}{Conceptual}{Analizar, Evaluar, Crear}

        \begin{document}

        \FichaCabecera

        \begin{materialesbox}
        \begin{itemize}
    \item Ficha impresa
    \item Lapiz
    \item Regla
    \item Calculadora sencilla si se necesita
\end{itemize}
        \end{materialesbox}

        \begin{pasosbox}
        \begin{enumerate}
    \item Lee la pregunta y decide primero si habla de borde o de interior.
    \item Anota la palabra clave: perimetro o area.
    \item Haz el calculo solo despues de elegir la magnitud correcta.
    \item Revisa la unidad final.
\end{enumerate}
        \end{pasosbox}

        \begin{recordatoriobox}
        \begin{itemize}
    \item Perimetro = borde exterior.
    \item Area = superficie interior.
    \item Perimetro usa m; area usa $m^2$.
\end{itemize}
        \end{recordatoriobox}

        \begin{apoyobox}
        \begin{itemize}
    \item Puedes dibujar una flecha al borde o sombrear el interior.
    \item Haz una tabla de dos columnas si te ayuda a separar ideas.
    \item Comprueba si el resultado deberia ser lineal o cuadrado.
\end{itemize}
        \end{apoyobox}

        \BloomSection{analizar}

\BloomExercise{analizar}{1}{Partes y relaciones}{Analiza un ejemplo relacionado con \emph{Perimetro y area: no son lo mismo}. Separa sus partes, indica cuales son relevantes y explica como se relacionan entre si.}{8}

\BloomExercise{analizar}{2}{Detecta errores o diferencias}{Compara dos respuestas, dos representaciones o dos procedimientos relacionados con esta sesion. Analiza sus diferencias y explica que errores o mejoras detectas.}{8}

\BloomSection{evaluar}

\BloomExercise{evaluar}{1}{Valora una propuesta}{Evalua una posible solucion o producto relacionado con \emph{Perimetro y area: no son lo mismo}. Decide si es adecuado y justifica tu opinion con criterios matematicos claros.}{8}

\BloomExercise{evaluar}{2}{Elige la mejor opcion}{Entre varias opciones posibles para cumplir este objetivo: \emph{Comparar perimetro y area para no confundir ambas medidas al analizar zonas urbanas.}, elige la mejor y argumenta por que. Incluye al menos dos criterios de valoracion.}{8}

\BloomSection{crear}

\BloomExercise{crear}{1}{Diseño propio}{Crea una produccion nueva relacionada con \emph{Perimetro y area: no son lo mismo}: un problema, una representacion, una mini guia, una explicacion visual o una propuesta de mejora.}{9}

\BloomExercise{crear}{2}{Alternativa mejorada}{Inventa una alternativa mas clara, mas util o mas creativa para trabajar este contenido. Debe estar alineada con el objetivo de la sesion y con la dimension conceptual.}{9}

        \begin{evidenciabox}
        \begin{itemize}
    \item Elegir la magnitud correcta antes de calcular.
    \item Comparar figuras con mismo perimetro o misma area.
    \item Justificar por que dos medidas distintas responden a preguntas distintas.
\end{itemize}
        \end{evidenciabox}

        \end{document}
